\documentclass[10pt,twocolumn]{article}

\usepackage[T1]{fontenc}
\usepackage{amsmath}
\usepackage{booktabs}
\usepackage[font=small,labelfont=bf]{caption}
\usepackage{graphicx}
\usepackage[margin=0.72in,columnsep=0.24in]{geometry}
\usepackage[colorlinks=true,allcolors=blue]{hyperref}
\usepackage{microtype}
\usepackage[numbers,sort&compress]{natbib}
\usepackage{placeins}
\usepackage{xcolor}

\graphicspath{{../figures/out/}}
\newcommand{\system}{\textsc{AutoTree}}
\newcommand{\fixturewarning}{\textbf{Fixture-harness illustration; not measured model performance.}\ }
\makeatletter
\setlength{\@fptop}{0pt}
\setlength{\@dblfptop}{0pt}
\makeatother

\title{AutoTree: Tree-Structured KV Execution for Test-Time Compute}
\author{Anonymous Authors}
\date{}

\begin{document}
\maketitle

\begin{abstract}
Reasoning workloads spend increasing amounts of inference compute exploring
candidate continuations, yet conventional serving treats each continuation as
an independent sequence.  We present \system, a system design and CPU-testable
vertical slice for executing reasoning as a tree.  Its Tree-KV substrate shares
paged key--value state across branches, forks with copy-on-write semantics,
deduplicates identical full pages, and restricts attention to each branch's
root-to-node path.  A deterministic Rust scheduler implements beam,
best-first, and Monte Carlo tree-search policies under hard token budgets, while
an OpenAI-shaped wire protocol exposes branch lifecycle events and reconciled
reuse counters.  This draft defines the intended GPU evaluation protocol and a
reproducible seven-figure publication pipeline.  It reports no GPU-scale or
large-model benchmark result: every plotted value is visibly marked as fixture
data and exists only to validate schemas, layout, error bars, and paper
integration.  Large-model quality, cost, and throughput claims remain pending.
\end{abstract}

\section{Introduction}

Test-time compute improves a model's opportunity to search, verify, and revise,
but it also changes the serving problem.  Independent best-of-$n$ sampling
repeats prompt and reasoning prefixes, retains branches that have already become
unlikely to help, and provides no shared budget controller.  The systems
question is therefore not only how to decode a sequence quickly, but how to
execute a dynamically changing \emph{tree} of sequences without losing the
memory and batching benefits of modern serving engines.

\system addresses this question with three contracts.  First, Tree-KV makes
branch ancestry explicit in the paged KV representation.  Second, a Rust policy
engine converts token and value events into fork, continue, kill, and finalize
commands under a global token budget.  Third, a versioned HTTP and streaming
surface exposes the winning answer together with enough branch and memory
counters to audit the execution.  These contracts allow a CPU reference path
to establish semantics before a production GPU runtime replaces it.

This paper makes four contributions:
\begin{itemize}
  \item a copy-on-write paged-KV tree model with full-page content deduplication
        and immediate reclamation after legal leaf pruning;
  \item branch-aware tree attention in which each query attends only to its
        root-to-node context;
  \item deterministic beam, best-first, and MCTS scheduling with explicit
        scorer synchronization and hard token-budget enforcement; and
  \item a provenance-checked figure pipeline and an evaluation protocol that
        refuse to turn fixture data into benchmark claims.
\end{itemize}

The present artifact is deliberately narrower than the target system.  It is a
CPU-testable vertical slice, not yet a production SGLang fork, a distributed GPU
runtime, or evidence for cost or throughput improvements.  The distinction
between implemented semantics and future performance work is maintained
throughout the paper.

\begin{figure*}[t]
  \centering
  \includegraphics[width=0.94\textwidth]{04_tree_topology.pdf}
  \caption{A branch-trace rendering slot: node color encodes a value estimate
  and red crosses mark pruned branches. \fixturewarning The topology is
  synthetic and is not a real AIME solve; a real trace will replace it without
  changing the renderer.}
  \label{fig:topology}
\end{figure*}

\section{Method}

\subsection{Tree-KV state and copy-on-write forks}

For each transformer layer, the KV pool stores fixed-size pages with shape
$[P,S,H_{kv},D]$, where $P$ is the number of physical pages and the default page
size $S$ is 16 tokens.  A branch records its parent, its root-to-node token
count, and an ordered block table of physical page identifiers.  Forking a
branch increments page references and copies no covered page.  When a child
writes to a shared page, the pool first copies that page and redirects only the
writer.  Thus branch creation is cheap while divergence pays the copy cost
lazily.

Pruning is defined only for leaves.  It decrements every covered page reference
and immediately returns zero-reference pages to the allocator.  Full frozen
pages may additionally enter a content-addressed deduplication pass: pages with
identical raw bytes are redirected to one physical page and duplicates are
freed.  The pool exports logical tokens, physical tokens, used pages, and the
reuse multiplier
\begin{equation}
  R_{KV}=\frac{\text{logical token references}}
                 {\text{physical tokens stored}}.
\end{equation}
The served wire contract requires a positive physical count and $R_{KV}\geq1$.

Convergent branch merging is a scheduler-level operation distinct from the
deduplication pass.  The wire protocol defines a terminal branch-merge event and
a live target branch.  The present engine does not yet drive those merge events;
it implements pool-level full-page deduplication.  This separation prevents the
paper from treating a storage primitive as an already complete semantic merge
algorithm.

\subsection{Branch-aware tree attention}

At a decode step, each active branch contributes one query.  Its block table
and context length select exactly the root-to-node K/V path.  A query cannot
attend to tokens owned only by a sibling, which is the tree analogue of a causal
mask.  The normative reference implementation uses PyTorch and fp32
accumulation.  An import-guarded Triton decode path is intended to match that
reference within dtype-specific tolerances; fused prefill remains future work.

\subsection{Rust scheduling policies}

The scheduler owns a branch tree whose root begins active.  Engine events report
sampled tokens, branch exhaustion, and optional external value scores.  Commands
authorize exactly one continuation step, fork an active branch, kill it with a
stable reason, or finalize it.  Beam search retains the highest-ranked live
frontier, best-first selects the strongest available branch, and MCTS balances
estimated value with exploration.  Deterministic random-number handling makes
policy traces replayable for a fixed configuration and seed.

External scoring is synchronized with decoding: once a branch requests a value,
it cannot consume another token until the corresponding score arrives.  This
prevents a fast branch from outrunning the policy decision that is supposed to
govern it.

\subsection{Budget controller}

Every sampled token consumes one unit of the total tree budget.  Per-branch and
global exhaustion produce explicit kill reasons, and the engine reclaims killed
leaf state immediately.  The budget is therefore a hard execution invariant,
not a post-hoc accounting field.  A future thinking-effort adapter can map a
model-level control to width, depth, scorer cadence, and token budget while
retaining this invariant.

\section{Implementation}

The repository is divided by contract.  The Python core contains a
device-agnostic paged KV pool, tree state, PyTorch reference attention, an
optional Triton decode dispatch, and a CPU model executor.  The Rust crate
implements branch state, the three scheduling policies, deterministic seeds,
and optional PyO3 bindings.  The serving package exposes OpenAI-shaped chat and
tree completion endpoints, server-sent branch events, validation, and
Prometheus metrics.  The typed Python SDK validates the event lineage and
exports rollout traces for RL-shaped consumers.

The tree endpoint accepts a policy, branch count, token budget, and optional
scorer.  Successful streams start each branch before referring to it, terminate
every non-winner by prune or merge, and finish with exactly one winner.  The
reported text follows the complete root-to-winner lineage.  Usage counts all
emitted branch tokens, while tree summaries reconcile branch counts, terminal
events, per-branch token spend, scores, and KV reuse.

The current implementation prioritizes contract integration on a CPU-only
Windows development path.  The Triton route is not validated by that gate, the
real Tree-KV engine does not yet emit semantic merge events, and serving is not
yet the production SGLang-based runtime described by the longer-term blueprint.

\subsection{Reproducible figure pipeline}

The publication figures are generated from versioned JSON with one command:
\begin{verbatim}
cd figures
uv run python make_figures.py --out out/
\end{verbatim}
The command emits all seven figures as PDF and SVG at 300 dpi plus a manifest of
input, source, and artifact SHA-256 hashes.  Provenance is mandatory.  Any input
whose provenance kind is \texttt{fixture} receives a visible
\texttt{FIXTURE DATA} watermark, and renderers reject absent measurements rather
than inventing them.  Tests cover schemas, axes, labels, watermarks, hashes, and
byte-for-byte deterministic regeneration.

\section{Evaluation Plan and Fixture Illustrations}

\textbf{No model-performance result is reported in this draft.}  The bundled
ThoughtBench tasks and every plotted value are synthetic contract fixtures.
They validate the end-to-end reporting path, but they cannot support a claim
about accuracy, cost, latency, KV reuse, or throughput.  GPU-scale measurements
and real benchmark traces are pending.

\subsection{Planned protocol}

The planned evaluation compares identical model weights and decoding settings
under stock serving baselines, sequential best-of-$n$, AutoTree-OFF, and
\system.  Initial task suites are AIME 2026, GPQA Diamond, and LiveCodeBench,
followed by broader coding and agentic suites.  Each configuration uses three
protocol seeds, fixed budgets and sampling parameters, automated graders, and
downloadable per-sample JSON.  The primary comparisons are iso-accuracy cost and
iso-budget accuracy; a result is not accepted if model, prompt set, decoding, or
hardware differs across the compared systems.

For task outcomes $c_{i,j}\in\{0,1\}$, accuracy@$k$ is the fraction of tasks
with at least one correct answer among the first $k$ samples.  Cost per correct
answer divides total priced input and output tokens by the number of correct
samples.  Rollout throughput is normalized as completed rollouts per hour per
GPU.  KV reuse uses the logical-to-physical multiplier defined in
Section~2; useful-token ratio remains unavailable unless the endpoint reports
it explicitly.

\begin{figure*}[t]
  \centering
  \includegraphics[width=0.92\textwidth]{01_accuracy_vs_tokens.pdf}
  \caption{Accuracy-versus-token layout for four planned model slots, comparing
  sequential execution with AutoTree and showing one-standard-deviation error
  bars over three synthetic seeds. \fixturewarning Model names label future
  evaluation slots; no named model produced these values.}
  \label{fig:scaling}
\end{figure*}

\begin{figure*}[t]
  \centering
  \includegraphics[width=0.92\textwidth]{02_cost_accuracy_pareto.pdf}
  \caption{Cost-per-correct versus accuracy Pareto layout for vLLM, SGLang, and
  AutoTree slots, with one-standard-deviation error bars over three synthetic
  seeds. \fixturewarning The curves are synthetic schema and layout checks, not
  measured serving results.}
  \label{fig:pareto}
\end{figure*}

\begin{figure}[t]
  \centering
  \includegraphics[width=\columnwidth]{03_kv_reuse_heatmap.pdf}
  \caption{Depth-by-branching KV-reuse layout, expressed as the
  logical/physical token multiplier rather than a percentage.
  \fixturewarning The cells are synthetic and not measured reuse.}
  \label{fig:reuse}
\end{figure}

\begin{figure}[t]
  \centering
  \includegraphics[width=\columnwidth]{05_branching_factor_ablation.pdf}
  \caption{Branching-factor ablation structure for
  $k\in\{1,2,4,8,16,32\}$. \fixturewarning The apparent trend is synthetic and
  must not be interpreted as an optimal branch count.}
  \label{fig:branching}
\end{figure}

\begin{figure}[t]
  \centering
  \includegraphics[width=\columnwidth]{06_rollout_throughput.pdf}
  \caption{Rollout-throughput comparison slot, normalized per GPU.
  \fixturewarning Bar heights and error bars are synthetic; no throughput claim
  follows from this panel.}
  \label{fig:throughput}
\end{figure}

\begin{figure}[t]
  \centering
  \includegraphics[width=\columnwidth]{07_thinking_effort_sweep.pdf}
  \caption{Thinking-effort sweep structure with AutoTree ON/OFF series.
  \fixturewarning This is a data slot for a later Inkling evaluation, not an
  Inkling measurement.}
  \label{fig:effort}
\end{figure}

\subsection{Acceptance criteria}

The final evaluation will report all seeds and configs, preserve failed and
missing cells, and include parity checks across batch sizes.  Accuracy parity
must be judged with uncertainty rather than a single point estimate.  Cost and
throughput claims require the same hardware accounting boundary, while reuse
claims must reconcile with logical and physical token counters.  Figure
captions will lose the fixture warning only when their complete provenance
points to released benchmark artifacts.

\FloatBarrier

\section{Related Work}

\paragraph{Paged KV serving.}
PagedAttention and vLLM apply virtual-memory ideas to KV-cache allocation and
continuous batching, substantially reducing fragmentation and duplicated
reservation~\cite{kwon2023vllm}.  \system adopts paged storage but makes branch
ancestry, copy-on-write forks, and tree accounting first-class.

\paragraph{Prefix reuse and structured programs.}
SGLang introduced RadixAttention to reuse prefixes across structured language
model programs and requests~\cite{zheng2024sglang}.  AutoTree's target serving
architecture is intentionally close to that substrate.  Its distinct emphasis
is a single dynamically scheduled reasoning tree with explicit prune, merge,
winner-lineage, and budget semantics.

\paragraph{Efficient attention.}
FlashAttention reduces attention memory traffic through tiled exact
computation~\cite{dao2022flashattention}.  This work is complementary: a
branch-aware mask determines which path is legal, while an efficient kernel
determines how that legal attention is executed.

\paragraph{Speculative decoding.}
Speculative decoding uses a cheaper draft process and a verification rule to
accelerate autoregressive generation without changing the target distribution
under its assumptions~\cite{leviathan2023speculative,chen2023speculative}.
AutoTree instead schedules multiple reasoning alternatives and may prune them
using values and budgets.  Speculative token generation could be composed with
each branch, but it does not replace tree-level memory sharing or policy.

\paragraph{Deliberate search.}
Tree of Thoughts frames language-model inference as search over intermediate
thoughts with evaluation and backtracking~\cite{yao2023tree}.  AutoTree focuses
on the systems substrate needed to execute such search without treating every
branch as an unrelated sequence.

\section{Limitations}

The artifact has five major limitations.  First, no real large-model or
GPU-scale benchmark has been run, so the paper cannot substantiate quality,
cost, or throughput improvements.  Second, the Windows verification path does
not validate Triton parity or performance.  Third, live semantic branch merging
is specified by the wire protocol but is not yet driven by the Tree-KV engine;
only full frozen-page deduplication exists in the storage layer.  Fourth, the
current server is a CPU vertical slice rather than the planned SGLang-based,
multi-GPU production runtime.  Fifth, value-guided pruning can discard a branch
that would later recover; scorer calibration and search bias therefore require
task-level study, not only systems benchmarks.

The paper's own presentation has a corresponding limitation: all seven figures
contain synthetic fixture values.  Visible watermarks, captions, provenance
checks, and manifest hashes reduce the risk of accidental misuse, but readers
should treat the panels only as executable templates for future measurements.

\section{Conclusion}

AutoTree recasts reasoning inference as execution of a shared-state tree.  Its
current vertical slice establishes contracts for copy-on-write paged KV,
branch-local attention, deterministic Rust scheduling, hard token budgets, and
auditable streaming summaries.  The accompanying publication pipeline turns
versioned JSON into seven deterministic, provenance-marked figures and an
integrated paper draft.  The remaining work is empirical and substantial:
validate kernels on GPUs, complete semantic merging and production serving, and
replace every fixture panel with released large-model measurements under the
stated protocol.


\bibliographystyle{plainnat}
\bibliography{refs}
\end{document}
